\section{AHSME 1951}
        \begin{enumerate}
        	\item (\textit{AHSME 1951}) The percent that $M$ is greater than $N$ is:


 $\textbf{(A) \ } \frac {100(M - N)}{M} \qquad \textbf{(B) \ } \frac {100(M - N)}{N} \qquad \textbf{(C) \ } \frac {M - N}{N} \qquad \textbf{(D) \ } \frac {M - N}{M} \qquad \textbf{(E) \ } \frac {100(M + N)}{N}$


	\item (\textit{AHSME 1951}) A rectangular field is half as wide as it is long and is completely enclosed by $x$ yards of fencing. The area in terms of $x$ is:


 $(\textbf{A})\ \frac{x^2}2 \qquad (\textbf{B})\ 2x^2 \qquad (\textbf{C})\ \frac{2x^2}9 \qquad (\textbf{D})\ \frac{x^2}{18} \qquad (\textbf{E})\ \frac{x^2}{72}$


	\item (\textit{AHSME 1951}) If the length of a diagonal of a square is $a + b$, then the area of the square is:


 $\textbf{(A)}\ (a+b)^{2}\qquad\textbf{(B)}\ \frac{1}{2}(a+b)^{2}\qquad\textbf{(C)}\ a^{2}+b^{2}$
$\textbf{(D)}\ \frac{1}{2}(a^{2}+b^{2})\qquad\textbf{(E)}\ \text{none of these}$


	\item (\textit{AHSME 1951}) A barn with a flat roof is rectangular in shape, $10$ yd. wide, $13$ yd. long and $5$ yd. high.  It is to be painted inside and outside, and on the ceiling, but not on the roof or floor. The total number of sq. yd. to be painted is:


 $\textbf{(A) \ } 360 \qquad \textbf{(B) \ } 460 \qquad \textbf{(C) \ } 490 \qquad \textbf{(D) \ } 590 \qquad \textbf{(E) \ } 720$


	\item (\textit{AHSME 1951}) Mr. A owns a home worth $\$10,000$.  He sells it to Mr. B at a $10 \%$ profit based on the worth of the house. Mr. B sells the house back to Mr. A at a $10 \%$ loss.  Then:


 $\textbf{(A)}\ \text{A comes out even} \qquad\textbf{(B)}\ \text{A makes \$1100 on the deal} \qquad\textbf{(C)}\ \text{A makes \$1000 on the deal}$
$\textbf{(D)}\ \text{A loses \$900 on the deal} \qquad\textbf{(E)}\ \text{A loses \$1000 on the deal}$


	\item (\textit{AHSME 1951}) The bottom, side, and front areas of a rectangular box are known. The product of these areas is equal to:


 $\textbf{(A)}\ \text{the volume of the box} \qquad\textbf{(B)}\ \text{the square root of the volume} \qquad\textbf{(C)}\ \text{twice the volume}$
$\textbf{(D)}\ \text{the square of the volume} \qquad\textbf{(E)}\ \text{the cube of the volume}$


	\item (\textit{AHSME 1951}) An error of $.02"$ is made in the measurement of a line $10"$ long, while an error of only $.2"$ is made in a measurement of a line $100"$ long. In comparison with the relative error of the first measurement, the relative error of the second measurement is:


 $\textbf{(A) \ } \text{greater by }.18 \qquad\textbf{(B) \ } \text{the same} \qquad \textbf{(C) \ } \text{less} \qquad\textbf{(D) \ } 10\text{ times as great} \qquad\textbf{(E) \ } \text{correctly described by both}$


	\item (\textit{AHSME 1951}) The price of an article is cut $10 \%.$ To restore it to its former value, the new price must be increased by:


 $\textbf{(A) \ } 10 \% \qquad\textbf{(B) \ } 9 \% \qquad \textbf{(C) \ } 11\frac{1}{9} \% \qquad\textbf{(D) \ } 11 \% \qquad\textbf{(E) \ } \text{none of these answers}$


	\item (\textit{AHSME 1951}) An equilateral triangle is drawn with a side length of $a.$ A new equilateral triangle is formed by joining the midpoints of the sides of the first one. then a third equilateral triangle is formed by joining the midpoints of the sides of the second; and so on forever. the limit of the sum of the perimeters of all the triangles thus drawn is:


 $\textbf{(A) \ } \text{Infinite} \qquad\textbf{(B) \ } 5\frac{1}{4}a \qquad \textbf{(C) \ } 2a \qquad\textbf{(D) \ } 6a \qquad\textbf{(E) \ } 4\frac{1}{2}a$


	\item (\textit{AHSME 1951}) Of the following statements, the one that is incorrect is:


 $\textbf{(A)}\ \text{Doubling the base of a given rectangle doubles the area.}\qquad\textbf{(B)}\ \text{Doubling the altitude of a triangle doubles the area.}$
$\textbf{(C)}\ \text{Doubling the radius of a given circle doubles the area.}$
$\textbf{(D)}\ \text{Doubling the divisor of a fraction and dividing its numerator by 2 changes the quotient.}$
$\textbf{(E)}\ \text{Doubling a given quantity may make it less than it originally was.}$


	\item (\textit{AHSME 1951}) The limit of the sum of an infinite number of terms in a geometric progression is $\frac {a}{1- r}$ where $a$ denotes the first term and $-1 < r < 1$ denotes the common ratio. The limit of the sum of their squares is:


 $\textbf{(A)}\ \frac {a^2}{(1 -r)^2} \qquad\textbf{(B)}\ \frac {a^2}{1 + r^2} \qquad\textbf{(C)}\ \frac {a^2}{1 - r^2} \qquad\textbf{(D)}\ \frac {4a^2}{1+ r^2} \qquad\textbf{(E)}\ \text{none of these}$


	\item (\textit{AHSME 1951}) At $2: 15$ o'clock, the hour and minute hands of a clock form an angle of:


 $\textbf{(A)}\ 30^{\circ} \qquad\textbf{(B)}\ 5^{\circ} \qquad\textbf{(C)}\ 22\frac {1}{2}^{\circ} \qquad\textbf{(D)}\ 7\frac {1}{2} ^{\circ} \qquad\textbf{(E)}\ 28^{\circ}$


	\item (\textit{AHSME 1951}) $A$ can do a piece of work in $9$ days. $B$ is $50\%$ more efficient than $A$. The number of days it takes $B$ to do the same piece of work is:


 $\textbf{(A)}\ 13\frac {1}{2} \qquad\textbf{(B)}\ 4\frac {1}{2} \qquad\textbf{(C)}\ 6 \qquad\textbf{(D)}\ 3 \qquad\textbf{(E)}\ \text{none of these answers}$


	\item (\textit{AHSME 1951}) In connection with proof in geometry, indicate which one of the following statements is \textit{incorrect}: 


 $\textbf{(A)}\ \text{Some statements are accepted without being proved.}$
$\textbf{(B)}\ \text{In some cases there is more than one correct order in proving certain propositions.}$
$\textbf{(C)}\ \text{Every term used in a proof must have been defined previously.}$
$\textbf{(D)}\ \text{It is not possible to arrive by correct reasoning at a true conclusion if, in the given, there is an untrue proposition.}$
$\textbf{(E)}\ \text{Indirect proof can be used whenever there are two or more contrary propositions.}$


	\item (\textit{AHSME 1951}) The largest number by which the expression $n^3-n$ is divisible for all possible integer values of $n$, is: 


 $\textbf{(A)}\ 2\qquad\textbf{(B)}\ 3\qquad\textbf{(C)}\ 4\qquad\textbf{(D)}\ 5\qquad\textbf{(E)}\ 6$


	\item (\textit{AHSME 1951}) If in applying the quadratic formula to a quadratic equation


 
 \[f(x) \equiv ax^2 + bx + c = 0,\]


 it happens that $c = \frac{b^2}{4a}$, then the graph of $y = f(x)$ will certainly:


 $\textbf{(A)}\ \text{have a maximum}\qquad\textbf{(B)}\ \text{have a minimum}\qquad\textbf{(C)}\ \text{be tangent to the x-axis}\\ \qquad\textbf{(D)}\ \text{be tangent to the y-axis}\qquad\textbf{(E)}\ \text{lie in one quadrant only}$


	\item (\textit{AHSME 1951}) Indicate in which one of the following equations $y$ is neither directly nor inversely proportional to $x$: 


 $\textbf{(A)}\ x+y = 0\qquad\textbf{(B)}\ 3xy = 10\qquad\textbf{(C)}\ x = 5y \qquad\textbf{(D)}\ 3x+y = 10 \qquad\textbf{(E)}\ x/y = \sqrt{3}\qquad$


	\item (\textit{AHSME 1951}) The expression $21x^2+ax+21$ is to be factored into prime binomial factors and without a numerical monomial factor. This can be done if the value ascribed to $a$ is:


 $\textbf{(A)}\ \text{any odd number}\qquad\textbf{(B)}\ \text{some odd number}\qquad\textbf{(C)}\ \text{any even number}\qquad\textbf{(D)}\ \text{some even number}\qquad\textbf{(E)}\ \text{zero}$


	\item (\textit{AHSME 1951}) A six place number is formed by repeating a three place number; for example, $256256$ or $678678$, etc. Any number of this form is always exactly divisible by: 


 $\textbf{(A)}\ 7\text{ only}\qquad\textbf{(B)}\ 11\text{ only}\qquad\textbf{(C)}\ 13\text{ only}\qquad\textbf{(D)}\ 101\qquad\textbf{(E)}\ 1001$


	\item (\textit{AHSME 1951}) When simplified and expressed with negative exponents, the expression $(x+y)^{-1}(x^{-1}+y^{-1})$ is equal to: 


 $\textbf{(A)}\ x^{-2}+2x^{-1}y^{-1}+y^{-2}\qquad\textbf{(B)}\ x^{-2}+2^{-1}x^{-1}y^{-1}+y^{-2}\qquad\textbf{(C)}\ x^{-1}y^{-1}$
$\textbf{(D)}\ \text{some even number}\qquad\textbf{(E)}\ \text{zero}$


	\item (\textit{AHSME 1951}) Given: $x > 0, y > 0, x > y$ and $z\not = 0$. The inequality which is not always correct is: 


 $\textbf{(A)}\ x+z > y+z\qquad\textbf{(B)}\ x-z > y-z\qquad\textbf{(C)}\ xz > yz$
$\textbf{(D)}\ \frac{x}{z^{2}}>\frac{y}{z^{2}}\qquad\textbf{(E)}\ xz^{2}> yz^{2}$


	\item (\textit{AHSME 1951}) The values of $a$ in the equation: $\log_{10}(a^{2}-15a) = 2$ are:


 $\textbf{(A)}\ \frac{15\pm\sqrt{233}}{2}\qquad\textbf{(B)}\ 20,-5\qquad\textbf{(C)}\ \frac{15\pm\sqrt{305}}{2}\qquad\textbf{(D)}\ \pm20$
$\textbf{(E)}\ \text{none of these}$


	\item (\textit{AHSME 1951}) The radius of a cylindrical box is $8$ inches and the height is $3$ inches. The number of inches that may be added to either the radius or the height to give the same nonzero increase in volume is: 


 $\textbf{(A)}\ 1\qquad\textbf{(B)}\ 5\frac{1}{3}\qquad\textbf{(C)}\ \text{any number}\qquad\textbf{(D)}\ \text{non-existent}\qquad\textbf{(E)}\ \text{none of these}$


	\item (\textit{AHSME 1951}) $\frac{2^{n+4}-2(2^{n})}{2(2^{n+3})}$ when simplified is:


 $\textbf{(A)}\ 2^{n+1}-\frac{1}{8}\qquad\textbf{(B)}\ -2^{n+1}\qquad\textbf{(C)}\ 1-2^{n}\qquad\textbf{(D)}\ \frac{7}{8}\qquad\textbf{(E)}\ \frac{7}{4}$


	\item (\textit{AHSME 1951}) The apothem of a square having its area numerically equal to its perimeter is compared with the apothem of an equilateral triangle having its area numerically equal to its perimeter. The first apothem will be: 


 $\textbf{(A)}\ \text{equal to the second}\qquad\textbf{(B)}\ \frac{4}{3}\text{ times the second}\qquad\textbf{(C)}\ \frac{2}{\sqrt{3}}\text{ times the second}\\ \textbf{(D)}\ \frac{\sqrt{2}}{\sqrt{3}}\text{ times the second}\qquad\textbf{(E)}\ \text{indeterminately related to the second}$


	\item (\textit{AHSME 1951}) In the equation $\frac{x(x-1)-(m+1)}{(x-1)(m-1)}=\frac{x}{m}$ the roots are equal when:


 $\textbf{(A)}\ m = 1\qquad\textbf{(B)}\ m =\frac{1}{2}\qquad\textbf{(C)}\ m = 0\qquad\textbf{(D)}\ m =-1\qquad\textbf{(E)}\ m =-\frac{1}{2}$


	\item (\textit{AHSME 1951}) Through a point inside a triangle, three lines are drawn from the vertices to the opposite sides forming six triangular sections. Then: 


 $\textbf{(A)}\ \text{the triangles are similar in opposite pairs}\qquad\textbf{(B)}\ \text{the triangles are congruent in opposite pairs}$
$\textbf{(C)}\ \text{the triangles are equal in area in opposite pairs}\qquad\textbf{(D)}\ \text{three similar quadrilaterals are formed}$
$\textbf{(E)}\ \text{none of the above relations are true}$


	\item (\textit{AHSME 1951}) The pressure $(P)$ of wind on a sail varies jointly as the area $(A)$ of the sail and the square of the velocity $(V)$ of the wind. The pressure on a square foot is $1$ pound when the velocity is $16$ miles per hour. The velocity of the wind when the pressure on a square yard is $36$ pounds is: 


 $\textbf{(A)}\ 10\frac{2}{3}\text{ mph}\qquad\textbf{(B)}\ 96\text{ mph}\qquad\textbf{(C)}\ 32\text{ mph}\qquad\textbf{(D)}\ 1\frac{2}{3}\text{ mph}\qquad\textbf{(E)}\ 16\text{ mph}$


	\item (\textit{AHSME 1951}) Of the following sets of data the only one that does not determine the shape of a triangle is: 


 $\textbf{(A)}\ \text{the ratio of two sides and the inc{}luded angle}\\ \qquad\textbf{(B)}\ \text{the ratios of the three altitudes}\\ \qquad\textbf{(C)}\ \text{the ratios of the three medians}\\ \qquad\textbf{(D)}\ \text{the ratio of the altitude to the corresponding base}\\ \qquad\textbf{(E)}\ \text{two angles}$


	\item (\textit{AHSME 1951}) If two poles $20''$ and $80''$ high are $100''$ apart, then the height of the intersection of the lines joining the top of each pole to the foot of the opposite pole is: 


 $\textbf{(A)}\ 50''\qquad\textbf{(B)}\ 40''\qquad\textbf{(C)}\ 16''\qquad\textbf{(D)}\ 60''\qquad\textbf{(E)}\ \text{none of these}$


	\item (\textit{AHSME 1951}) A total of $28$ handshakes were exchanged at the conclusion of a party. Assuming that each participant was equally polite toward all the others, the number of people present was: 


 $\textbf{(A)}\ 14\qquad\textbf{(B)}\ 28\qquad\textbf{(C)}\ 56\qquad\textbf{(D)}\ 8\qquad\textbf{(E)}\ 7$


	\item (\textit{AHSME 1951}) If $\triangle ABC$ is inscribed in a semicircle whose diameter is $AB$, then $AC+BC$ must be 


 $\textbf{(A)}\ \text{equal to }AB\qquad\textbf{(B)}\ \text{equal to }AB\sqrt{2}\qquad\textbf{(C)}\ \geq AB\sqrt{2}\qquad\textbf{(D)}\ \leq AB\sqrt{2} \qquad \textbf{(E)}\ AB^{2}$


	\item (\textit{AHSME 1951}) The roots of the equation $x^{2}-2x = 0$ can be obtained graphically by finding the abscissas of the points of intersection of each of the following pairs of equations except the pair: 


 $\textbf{(A)}\ y = x^{2}, y = 2x\qquad\textbf{(B)}\ y = x^{2}-2x, y = 0\qquad\textbf{(C)}\ y = x, y = x-2\qquad\textbf{(D)}\ y = x^{2}-2x+1, y = 1$
$\textbf{(E)}\ y = x^{2}-1, y = 2x-1$


 \textit{[Note: Abscissa means x-coordinate.]}


 

	\item (\textit{AHSME 1951}) The value of $10^{\log_{10}7}$ is: 


 $\textbf{(A)}\ 7\qquad\textbf{(B)}\ 1\qquad\textbf{(C)}\ 10\qquad\textbf{(D)}\ \log_{10}7\qquad\textbf{(E)}\ \log_{7}10$


	\item (\textit{AHSME 1951}) If $a^{x}= c^{q}= b$ and $c^{y}= a^{z}= d$, then


 $\textbf{(A)}\ xy = qz\qquad\textbf{(B)}\ \frac{x}{y}=\frac{q}{z}\qquad\textbf{(C)}\ x+y = q+z\qquad\textbf{(D)}\ x-y = q-z$
$\textbf{(E)}\ x^{y}= q^{z}$


	\item (\textit{AHSME 1951}) Which of the following methods of proving a geometric figure a locus is not correct? 


 $\textbf{(A)}\ \text{Every point of the locus satisfies the conditions and every point not on the locus does not satisfy the conditions.}$
$\textbf{(B)}\ \text{Every point not satisfying the conditions is not on the locus and every point on the locus does satisfy the conditions.}$
$\textbf{(C)}\ \text{Every point satisfying the conditions is on the locus and every point on the locus satisfies the conditions.}$
$\textbf{(D)}\ \text{Every point not on the locus does not satisfy the conditions and every point not satisfying}\\ \text{the conditions is not on the locus.}$
$\textbf{(E)}\ \text{Every point satisfying the conditions is on the locus and every point not satisfying the conditions is not on the locus.}$


	\item (\textit{AHSME 1951}) A number which when divided by $10$ leaves a remainder of $9$, when divided by $9$ leaves a remainder of $8$, by $8$ leaves a remainder of $7$, etc., down to where, when divided by $2$, it leaves a remainder of $1$, is: 


 $\textbf{(A)}\ 59\qquad\textbf{(B)}\ 419\qquad\textbf{(C)}\ 1259\qquad\textbf{(D)}\ 2519\qquad\textbf{(E)}\ \text{none of these answers}$


	\item (\textit{AHSME 1951}) A rise of $600$ feet is required to get a railroad line over a mountain. The grade can be kept down by lengthening the track and curving it around the mountain peak. The additional length of track required to reduce the grade from $3\%$ to $2\%$ is approximately: 


 $\textbf{(A)}\ 10000\text{ ft.}\qquad\textbf{(B)}\ 20000\text{ ft.}\qquad\textbf{(C)}\ 30000\text{ ft.}\qquad\textbf{(D)}\ 12000\text{ ft.}\qquad\textbf{(E)}\ \text{none of these}$


	\item (\textit{AHSME 1951}) A stone is dropped into a well and the report of the stone striking the bottom is heard $7.7$ seconds after it is dropped. Assume that the stone falls $16t^2$ feet in t seconds and that the velocity of sound is $1120$ feet per second. The depth of the well is:


 $\textbf{(A)}\ 784\text{ ft.}\qquad\textbf{(B)}\ 342\text{ ft.}\qquad\textbf{(C)}\ 1568\text{ ft.}\qquad\textbf{(D)}\ 156.8\text{ ft.}\qquad\textbf{(E)}\ \text{none of these}$


	\item (\textit{AHSME 1951}) $\left(\frac{(x+1)^{2}(x^{2}-x+1)^{2}}{(x^{3}+1)^{2}}\right)^{2}\cdot\left(\frac{(x-1)^{2}(x^{2}+x+1)^{2}}{(x^{3}-1)^{2}}\right)^{2}$ equals:


 $\textbf{(A)}\ (x+1)^{4}\qquad\textbf{(B)}\ (x^{3}+1)^{4}\qquad\textbf{(C)}\ 1\qquad\textbf{(D)}\ [(x^{3}+1)(x^{3}-1)]^{2}$
$\textbf{(E)}\ [(x^{3}-1)^{2}]^{2}$


	\item (\textit{AHSME 1951}) The formula expressing the relationship between $x$ and $y$ in the table is:

 \[\begin{tabular}{|c|c|c|c|c|c|}\hline x & 2 & 3 & 4 & 5 & 6\\ \hline y & 0 & 2 & 6 & 12 & 20\\ \hline\end{tabular}\]
$\textbf{(A)}\ y = 2x-4\qquad\textbf{(B)}\ y = x^{2}-3x+2\qquad\textbf{(C)}\ y = x^{3}-3x^{2}+2x$
$\textbf{(D)}\ y = x^{2}-4x\qquad\textbf{(E)}\ y = x^{2}-4$


	\item (\textit{AHSME 1951}) If $x =\sqrt{1+\sqrt{1+\sqrt{1+\sqrt{1+\cdots}}}}$, then:


 $\textbf{(A)}\ x = 1\qquad\textbf{(B)}\ 0 < x < 1\qquad\textbf{(C)}\ 1 < x < 2\qquad\textbf{(D)}\ x\text{ is infinite}$
$\textbf{(E)}\ x > 2\text{ but finite}$


	\item (\textit{AHSME 1951}) Of the following statements, the only one that is incorrect is:


 $\textbf{(A)}\ \text{An inequality will remain true after each side is increased,}$ $\text{ decreased, multiplied or divided (zero excluded) by the same positive quantity.}$
$\textbf{(B)}\ \text{The arithmetic mean of two unequal positive quantities is greater than their geometric mean.}$
$\textbf{(C)}\ \text{If the sum of two positive quantities is given, their product is largest when they are equal.}$
$\textbf{(D)}\ \text{If }a\text{ and }b\text{ are positive and unequal, }\frac{1}{2}(a^{2}+b^{2})\text{ is greater than }[\frac{1}{2}(a+b)]^{2}.$
$\textbf{(E)}\ \text{If the product of two positive quantities is given, their sum is greatest when they are equal.}$


	\item (\textit{AHSME 1951}) If $\frac{xy}{x+y}= a,\frac{xz}{x+z}= b,\frac{yz}{y+z}= c$, where $a, b, c$ are other than zero, then $x$ equals: 


 $\textbf{(A)}\ \frac{abc}{ab+ac+bc}\qquad\textbf{(B)}\ \frac{2abc}{ab+bc+ac}\qquad\textbf{(C)}\ \frac{2abc}{ab+ac-bc}$
$\textbf{(D)}\ \frac{2abc}{ab+bc-ac}\qquad\textbf{(E)}\ \frac{2abc}{ac+bc-ab}$


	\item (\textit{AHSME 1951}) If you are given $\log 8\approx .9031$ and $\log 9\approx .9542$, then the only logarithm that cannot be found without the use of tables is:


 $\textbf{(A)}\ \log 17\qquad\textbf{(B)}\ \log\frac{5}{4}\qquad\textbf{(C)}\ \log 15\qquad\textbf{(D)}\ \log 600\qquad\textbf{(E)}\ \log .4$


	\item (\textit{AHSME 1951}) $AB$ is a fixed diameter of a circle whose center is $O$. From $C$, any point on the circle, a chord $CD$ is drawn perpendicular to $AB$. Then, as $C$ moves over a semicircle, the bisector of angle $OCD$ cuts the circle in a point that always: 


 $\textbf{(A)}\ \text{bisects the arc }AB\qquad\textbf{(B)}\ \text{trisects the arc }AB\qquad\textbf{(C)}\ \text{varies}$
$\textbf{(D)}\ \text{is as far from }AB\text{ as from }D\qquad\textbf{(E)}\ \text{is equidistant from }B\text{ and }C$


	\item (\textit{AHSME 1951}) If $r$ and $s$ are the roots of the equation $ax^2+bx+c=0$, the value of $\frac{1}{r^{2}}+\frac{1}{s^{2}}$ is:


 $\textbf{(A)}\ b^{2}-4ac\qquad\textbf{(B)}\ \frac{b^{2}-4ac}{2a}\qquad\textbf{(C)}\ \frac{b^{2}-4ac}{c^{2}}\qquad\textbf{(D)}\ \frac{b^{2}-2ac}{c^{2}}$
$\textbf{(E)}\ \text{none of these}$


	\item (\textit{AHSME 1951}) The area of a square inscribed in a semicircle is to the area of the square inscribed in the entire circle as: 


 $\textbf{(A)}\ 1: 2\qquad\textbf{(B)}\ 2: 3\qquad\textbf{(C)}\ 2: 5\qquad\textbf{(D)}\ 3: 4\qquad\textbf{(E)}\ 3: 5$


	\item (\textit{AHSME 1951}) The medians of a right triangle which are drawn from the vertices of the acute angles are $5$ and $\sqrt{40}$. The value of the hypotenuse is:


 $\textbf{(A)}\ 10\qquad\textbf{(B)}\ 2\sqrt{40}\qquad\textbf{(C)}\ \sqrt{13}\qquad\textbf{(D)}\ 2\sqrt{13}\qquad\textbf{(E)}\ \text{none of these}$


	\item (\textit{AHSME 1951}) Tom, Dick and Harry started out on a $100$-mile journey. Tom and Harry went by automobile at the rate of $25$ mph, while Dick walked at the rate of $5$ mph. After a certain distance, Harry got off and walked on at $5$ mph, while Tom went back for Dick and got him to the destination at the same time that Harry arrived. The number of hours required for the trip was: 


 $\textbf{(A)}\ 5\qquad\textbf{(B)}\ 6\qquad\textbf{(C)}\ 7\qquad\textbf{(D)}\ 8\qquad\textbf{(E)}\ \text{none of these answers}$


\end{enumerate}